% Root derivation. HZA: full carrier, zero graph, and verified calculation scope.
\section{The aligned Hurwitz zero graph with its complete support map}
\label{sec:hurwitz-alignment-full-map}

\paragraph{Sources and actual objects.}
The original programme source is the retained April 2026 draft
\emph{Split-Zero Support, Residue-Parity Reversal, Shifted-Hurwitz CUE
Defects, and Finite Phase-Space Enlargement}, printed byline
``Compiled theorem draft'', source labels \texttt{thm:universal-flow},
\texttt{prop:non-euler}, and \texttt{thm:zero-motion}.
Its complete supplied TeX is retained with this calculation; its human
author is not specified and is not inferred here. The following proof
keeps its original $s,t$ coordinates. The operator construction in HPS
uses the arithmetic scaling of Jean-Beno\^it Bost and Alain Connes,
and the precise endomotive presentation of Alain Connes, Caterina
Consani and Matilde Marcolli,
\href{https://arxiv.org/abs/0806.2401v1}{\emph{Fun with $\mathbb F_1$}},
source equations \texttt{sigmatgenerators} and relations (c1)--(c4).
The full split carrier is the one in
\href{https://github.com/KokunoYumeto/zeta-function-research-reader/blob/a2851f9eb352e7e4376bc629de86fa4ed9c1c434/workbenches/splitzero-tandem/foundations/split-support-geometry-v11/split_support_geometry_arithmetic_curve_v11.tex}{\emph{Split Support Geometry}, v11, Definition \texttt{def:lattice-split}}.
The new claims below are proved here; numerical path labels have the
explicit, narrower meaning stated at the end.

\subsection{The full carrier and its finite-cardinality receiver}

Let $L$ be the original nontrivial bounded distributive lattice, and $R$
a commutative ring. Retain, without identifying any labelled zeros,
\begin{equation}\tag{HZA1}
\begin{gathered}
G_L(R)=\{(0,\lambda):\lambda\in L\}\cup(R\times\{1_L\}),\qquad
\tau=(0,0_L),\quad e=(0,1_L),\\
(a,\lambda)+(b,\mu)=(a+b,\lambda\vee\mu),\qquad
(a,\lambda)(b,\mu)=(ab,\lambda\wedge\mu),\\
K_L=K\otimes_{\mathbb B}L,\qquad
\iota_L(\lambda)=1_K\otimes\lambda,\qquad
\Phi_R(a,\lambda)=(a,\iota_L(\lambda)).
\end{gathered}
\end{equation}
Here $K$ is the programme's tropical semifield, and $\mathbb B$ its
Boolean subsemiring. The map $K\to\mathbb B$ sending zero to zero and
every nonzero element to one preserves addition and multiplication.
Tensoring it with $L$ gives $\rho_L:K_L\to L$ with
$\rho_L\iota_L=\operatorname{id}_L$. Thus $\Phi_R$ is injective,
preserves both operations and recovers every label, including labels
whose amplitude cancels. This supplies the actual comparison map,
rather than a count of extra points.

For finite $L$ and $n\ge1$, the exact underlying-set count is
\begin{equation}\tag{HZA2}
\left|G_L(\mathbb Z/n\mathbb Z)\right|=(n-1)+|L|,
\qquad
\sum_{n=1}^{\infty}\left|G_L(\mathbb Z/n\mathbb Z)\right|^{-s}
=\zeta(s,|L|),\quad \operatorname{Re}s>1.
\end{equation}
Indeed there are precisely $n-1$ nonzero amplitudes, each forced to
have top label, and precisely $|L|$ distinct zero-amplitude elements.
The sets are disjoint. Substituting $n-1+|L|$ in the absolutely
convergent sum proves the second equality. The quotient
$(a,\lambda)\mapsto a$ identifies all these zeros and gives the
ordinary count $n$. For $|L|=2$, the separated count gives
$\zeta(s,2)=\zeta(s)-1$. For arbitrary finite $L$ its shift is
$t=|L|-1$, with no change in the carrier (HZA1).

The count is only this specified receiver. It does not recover the
meet, join, or prime spectrum of $L$. For example, the four-element
chain and the four-element Boolean lattice have the same count;
they are not isomorphic because one is totally ordered and the other
has two incomparable elements. For infinite $L$ the displayed finite
cardinality series is not defined. The full map (HZA1) remains
defined. The real continuous shift below is an analytic family;
it is not an assertion that the cardinality of $L$ varies continuously.

\subsection{An analytic construction through the critical strip}

Put
\begin{equation}\tag{HZA3}
 F(s,t)=\zeta(s,1+t),\qquad
 h_t(x)=\frac{x e^{-(1+t)x}}{1-e^{-x}}.
\end{equation}
The branch of $\log(1+t)$ and of $\log(n+t)$ for complex $t$ is the
one on $\operatorname{Re}(1+t)>0$. On
$\operatorname{Re}s>0$, $s\ne1$, $\operatorname{Re}(1+t)>0$, the
following absolutely convergent expression constructs $F$:
\begin{equation}\tag{HZA4}
 F(s,t)=\frac1{\Gamma(s)}\left\{
 \frac1{s-1}+
 \int_0^1 x^{s-2}(h_t(x)-1)\,dx+
 \int_1^\infty x^{s-2}h_t(x)\,dx\right\}.
\end{equation}
Near zero, $h_t(x)=1+O(x)$, uniformly on compact parameter sets.
For large positive $x$, it is bounded by a constant times
$x e^{-c x}$, where $c>0$ is a lower bound for
$\operatorname{Re}(1+t)$ on the chosen compact set. The same bounds
after any fixed number of $s,t$ derivatives acquire only powers of
$x$ and $|\log x|$. These bounds prove absolute convergence and local
uniform differentiability under both integrals. For
$\operatorname{Re}s>1$, combine the two integrals with
$\int_0^1x^{s-2}\,dx=(s-1)^{-1}$ and expand
$(1-e^{-x})^{-1}=\sum_{j\ge0}e^{-jx}$. Absolute convergence gives
$\Gamma(s)\sum_{n\ge1}(n+t)^{-s}$. This proves (HZA4) in the original
half-plane and supplies its holomorphic continuation to the stated
domain. Its only possible pole there is $s=1$, with residue one.

Differentiating $h_t$ gives $\partial_t h_t=-x h_t$, so the two
integrals immediately give the exact flow equation
\begin{equation}\tag{HZA5}
 \partial_t F(s,t)=-s F(s+1,t),\qquad
 \partial_t^kF(s,t)=(-1)^k(s)_kF(s+k,t).
\end{equation}
The first equality follows from (HZA4) and the convergent Mellin
integral for $F(s+1,t)$; the second follows by induction, with
$(s)_k=s(s+1)\cdots(s+k-1)$ and $(s)_0=1$. Holomorphy in $|t|<1$
then gives the actual, locally uniformly convergent Taylor series
\begin{equation}\tag{HZA6}
 F(s,t)=\sum_{k=0}^{\infty}
       \frac{(-t)^k(s)_k}{k!}\zeta(s+k),\quad
       \operatorname{Re}s>0,\ s\ne1,\ |t|<1.
\end{equation}
In this domain the factors on the right introduce no additional pole:
for $k\ge1$ the argument has real part greater than one.
For any compact subset in $s$ and any smaller closed $t$ disc,
Cauchy's integral formula in a larger disc of radius less than one
proves the claimed uniform convergence. Thus no formal exponential
is being used as an unproved analytic operator assertion.

\subsection{Every nontrivial zero gives one exact local graph}

Let $\rho$ be any zero of $\zeta$ in $0<\operatorname{Re}\rho<1$,
and let $m$ be its order. This is a quantification over actual zeros;
no zero on or off a particular vertical line is supplied as an
assumption about all the zeros. The Euler product in the absolutely
convergent half-plane gives
\begin{equation}\tag{HZA7}
 B_\rho=\rho\zeta(\rho+1)\ne0,
 \qquad \partial_t F(\rho,0)=-B_\rho.
\end{equation}
For completeness, nonvanishing of that Euler product follows from
$\sum_p\sum_{k\ge1}|p^{-k(\rho+1)}|/k<\infty$: its exponential is
nonzero, and absolute expansion by unique prime factorization equals
the zeta series. Also $\rho\ne0$ in the stated strip.

There is a neighbourhood of $\rho$ and a unique holomorphic function
$\theta$ there, taking values near zero, for which the full local
zero set is
\begin{equation}\tag{HZA8}
 \{(s,t):F(s,t)=0\}=\{(s,\theta(s))\},\qquad \theta(\rho)=0.
\end{equation}
Here is a construction with the uniqueness included. Write
$B=\partial_t F(\rho,0)\ne0$. Choose sufficiently small closed
discs $|s-\rho|\le d$, $|t|\le r$ inside the holomorphy domain so
that $|1-B^{-1}\partial_tF(s,t)|\le1/4$ and
$|B^{-1}F(s,0)|\le r/2$. Such discs exist by continuity and
$F(\rho,0)=0$. The map $t\mapsto t-B^{-1}F(s,t)$ sends the
$t$ disc into itself and has Lipschitz constant at most $1/4$.
Starting from zero, its iterates converge uniformly to the unique
fixed point. Each iterate is holomorphic in $s$ on the open disc;
uniform convergence makes the limit holomorphic. A zero of $F$
in the product of discs is exactly a fixed point, proving (HZA8).

The order and leading coefficient of this graph are determined
without any simplicity assumption:
\begin{equation}\tag{HZA9}
 \theta(s)=(s-\rho)^m U_\rho(s),\qquad
 U_\rho(\rho)=\frac{\zeta^{(m)}(\rho)}{m!\rho\zeta(\rho+1)}\ne0.
\end{equation}
To prove this, the fundamental theorem of calculus gives
$F(s,t)=\zeta(s)+t A(s,t)$, where
$A(s,t)=\int_0^1\partial_tF(s,ut)\,du$ is holomorphic and
$A(\rho,0)=-B_\rho$. At the graph,
$\theta(s)=-\zeta(s)/A(s,\theta(s))$. Its denominator is a
holomorphic unit near $\rho$. Factoring the zero of $\zeta$ proves
both the order and the coefficient in (HZA9). Multiple zeta zeros
therefore give ramification of the projection to $t$, rather than a
failure to construct the total zero graph.

At a simple zero the graph can instead be written $s=\rho(t)$,
by the same contraction argument with the variables interchanged.
Implicit differentiation in (HZA5) proves
\begin{equation}\tag{HZA10}
 \rho'(0)=\frac{\rho\zeta(\rho+1)}{\zeta'(\rho)}.
\end{equation}
If this actual zero has $\operatorname{Re}\rho=1/2$ and the real
part of (HZA10) is nonzero, then
$\operatorname{Re}\rho(t)-1/2=t\operatorname{Re}\rho'(0)+O(t^2)$
for real $t$. Dividing by $t$ and using continuity proves a
transverse crossing for all sufficiently small nonzero $t$.
This conclusion is local. It places no restriction on a later
return of the same branch to the line.

\subsection{The supported graph and all its comparison maps}

For any open product of discs $D$ used above, let $R=\mathcal O(D)$.
The original function is the element $(F,1_L)$ of $G_L(R)$, and
its value at a point $x=(s,t)$ is obtained by the unital semiring
homomorphism
\begin{equation}\tag{HZA11}
\begin{array}{ccc}
 G_L(\mathcal O(D))&\xrightarrow{\Phi_{\mathcal O(D)}}&
                 \mathcal O(D)\times K_L\\
 {\scriptstyle(a,\lambda)\mapsto(a(x),\lambda)}\downarrow&&
       \downarrow{\scriptstyle(a,w)\mapsto(a(x),w)}\\
 G_L(\mathbb C)&\xrightarrow{\Phi_{\mathbb C}}&\mathbb C\times K_L.
\end{array}
\end{equation}
Evaluation preserves ring addition and multiplication, and the label
is unchanged, proving the left map preserves the displayed split
operations. Both composites are $(a(x),\iota_L\lambda)$, proving
commutativity on every element, including every labelled zero.
Consequently the zero graph just constructed is exactly
\begin{equation}\tag{HZA12}
 \{x:\operatorname{ev}_x(F,1_L)=e\}
 =\{x:\operatorname{ev}_x\Phi(F,1_L)=(0,1_{K_L})\}.
\end{equation}
Its value is supported zero, not $\tau$. All other zero sections
$(0,\lambda)$ remain present, evaluate to themselves, and are
recovered from the full map by $\rho_L$. The faithful comparison
does not impose an additional equation on $s$. The disc-algebra
prime operator in HPS gives the entire family another exact receiver;
HCR gives its residue-class receivers and their retained support
defects. These are calculations on the same $F$, with the displayed
maps, rather than a change of its zero convention.

\subsection{The actual bounded calculation and its limits}

The existing programme velocity certificate stores enclosures for
$V_j=2\operatorname{Re}\rho_j'(0)$. Its first 60 rows have 29
positive and 31 negative signs. The factor two must be retained when
comparing that file with (HZA10). The new script
\texttt{compute\_hurwitz\_paths.py} independently recomputes those
60 velocities at 45 decimal digits, checks each sign, and solves
for actual nearby Hurwitz zeros at $t=10^{-4}$ and $t=-10^{-4}$.
It checks residuals less than $10^{-33}$. These are numerical checks;
the pre-existing interval enclosures and the new non-interval
recomputations are recorded as different evidence.

The same script numerically follows the first five starting zeros
to $t=-1/2$ and $t=1$, with adaptive predictor/Newton steps and
80-digit endpoint refinement. At $t=1$, the first three endpoints
are, to the displayed decimal precision,
\[
 1.40778804065+23.32798775456i,\qquad
 0.41092175383+31.71541044219i,\qquad
 1.36768818112+38.41910556064i.
\]
The exact endpoint identities are
\begin{equation}\tag{HZA13}
 F(s,1)=\zeta(s)-1,\qquad
 F(s,-1/2)=(2^s-1)\zeta(s).
\end{equation}
The first follows by deleting the first term of the zeta series.
For the second, split that series into odd and even integers:
$\sum_{n\ge0}(n+1/2)^{-s}=2^s\sum_{n\text{ odd}}n^{-s}
=2^s(1-2^{-s})\zeta(s)$. Both identities extend from their
convergent half-plane by holomorphic continuation. The nonzero
factor $2^s$ is retained. The zeros of $2^s-1$ have
$s=2\pi i k/\log2$ for every integer $k$, including $k=0$.

The five paths in the figure are interpolations of numerical sample
points. They were initialized at known Riemann zeros; their common
value $\operatorname{Re}s=1/2$ at $t=0$ is therefore not an
independent proof for untracked zeros. The first and third numerical
paths reach the first two positive imaginary zeros of $2^s-1$;
the other three reach earlier Riemann zeros. No rigorous global
branch-identity certificate is claimed for this numerical path
tracking. Separate interval proofs in HXC certify two additional
transverse critical-line crossings at strictly positive shifts.
They show why the local conclusion (HZA10) must not be changed into
the assertion that a zero can cross only at an Euler-product parameter.
